\begin{scanpage}{0127}{110}
Now let us find the field scattered by the spherical segment considering the discontinuity of the surface; one may neglect the perturbation of the current as a consequence of the smooth curve of the surface if $ka\gg1$ [74]. The nonuniform part of the current which is caused by the discontinuity creates in the direction $\vartheta=\pi$ the field (17.03). Summing the latter with the field (19.08), we find the desired field
\[E_x=-H_y=-\frac{aE_{0x}}2\left(\frac1{\cos\omega}+\frac{\frac2n\sin\frac\pi n}{\cos\frac\pi n-\cos\frac{2\omega}n}e^{2ikl}\right)\frac{e^{ikR}}R.\tag{19.11}\]
Consequently, the effective scattering area of a spherical segment will equal
\[\begin{gathered}\sigma=\pi a^2\left|\frac1{\cos\omega}+\frac{\frac2n\sin\frac\pi n}{\cos\frac\pi n-\cos\frac{2\omega}n}e^{2ikl}\right|^2,\\n=1+\frac{\omega+\Omega}\pi.
\end{gathered}\tag{19.12}\]
In the physical optics approach, a similar quantity is determined by field (19.08) and equals
\[\sigma^0=\pi a^2\left|\frac1{\cos\omega}-\operatorname{tg}\omega\,e^{2ikl}\right|^2.\tag{19.13}\]
With the deforming of the spherical surface into a disk $(\omega\to\frac\pi2,\ l\to0,\ \Omega=\mathrm{const})$, Equations (19.12) and (19.13) are transformed, respectively, to the form
\[\left.\begin{aligned}\sigma&=\pi a^2\left|ika+\frac1n\operatorname{ctg}\frac\pi n\right|^2,\\\sigma^0&=\pi a^2(ka)^2.\end{aligned}\right\}\tag{19.14}\]
It follows from Equations (19.12) and (19.13) that the effective scattering area of a spherical segment is an oscillating function of its length. The oscillation period equals $\frac\lambda2$. Numerical calculations performed on the basis of these equations showed (Figure 43) that, with small angles of the discontinuity $(\Omega=15^\circ)$, one may still neglect the field from the nonuniform part of the current. In Figure 44, graphs are constructed for the effective scattering area of a
\end{scanpage}
